\section{Plugin-Integration}

Die Erweiterbarkeit von Pablo wurde softwareseitig über eine klar definierte Plugin-Schnittstelle umgesetzt. 
Ziel war es, Pablo nicht auf einen festen Funktionsumfang zu beschränken, sondern zusätzliche Werkzeuge unabhängig vom Kernsystem ergänzen zu können.

\subsection{Ergebnis}

Die Plugins werden beim Start dynamisch aus dem Ordner \texttt{plugins} geladen. Berücksichtigt werden nur Module, die die beiden Bestandteile \texttt{TOOL\_DEF} und \texttt{execute} bereitstellen. Eine zusätzliche Registrierung im Kernsystem ist dadurch nicht erforderlich.
Jedes Plugin ist als eigenständiges Python-Modul aufgebaut. \texttt{TOOL\_DEF} beschreibt die für das Sprachmodell relevante Schnittstelle mit Name, Beschreibung und Parametern, während \texttt{execute(args)} die eigentliche Ausführung übernimmt und den Antworttext erzeugt. 
Zusätzliche Hilfsfunktionen, etwa für API-Aufrufe oder lokale Dateiverarbeitung, können innerhalb des jeweiligen Moduls ergänzt werden.
Neue Plugins können nach demselben Muster ergänzt und beim nächsten Systemstart automatisch geladen werden.
Im Prototyp wurden vier Basis-Plugins umgesetzt: Wetter, Notizen, Kalender und Timer. Damit kann Pablo neben freier Sprachverarbeitung auch konkrete unterstützende Funktionen im Alltag ausführen.

\subsection{Einordnung}

Die Trennung zwischen Basissystem und Plugins ermöglicht es, den Funktionsumfang von Pablo zu erweitern, ohne die zentrale Sprachverarbeitung oder den Hauptablauf des Systems grundlegend anzupassen. Im Entwicklungsprozess wurde diese modulare Struktur einer monolithischen Implementierung vorgezogen, da neue Funktionen so als eigenständige Module ergänzt werden können.
Die Plugin-Architektur unterstützt damit den Anspruch, Pablo als anpassbaren und erweiterbaren Companion zu entwickeln. Die im Prototyp umgesetzten Plugins decken typische Alltagsszenarien ab und bilden einen ersten funktionalen Kern, der sowohl die praktische Nutzbarkeit als auch die Umsetzbarkeit des modularen Erweiterungskonzepts verdeutlicht.